% ============================================================
% Chapter 2: State of the Art
%   Section: Automated Fact-Checking
%     Subsection: Deciding What to Search
% Included from chapters/02-state-of-the-art/01-automated-fact-checking/automated-fact-checking.tex
% Source: docs/troubleshooting/fact-checking/01-what-to-look-for.md (6 Nov);
% each paper checked against its ACL Anthology or arXiv page on 6 Nov 2026.
% ============================================================

\subsection{Deciding What to Search}
\label{subsec:sota-what-to-search}

Open retrieval needs a query, and the claim's sentence is rarely a good
one. A human fact-checker first decides what the claim turns on --- a
figure, a date, who did what --- and which source would settle it, and
only then searches. The literature has made that step explicit in three
ways.

\paragraph{Asking questions.} Fan et al.\ generate \emph{question-answering
briefs}: questions that decompose the claim, answered from the web
before a person checks it. Briefs raised crowdworkers' accuracy by about
10\% and cut volunteer fact-checkers' time by about a fifth
\cite{fan2020qabriefs}. Chen et al.\ distinguish \emph{literal}
subquestions, which ask about what the claim says, from \emph{implied}
ones, which ask about what it leaves unsaid: their best generator
recovered over 70\% of the literal questions and 18\% of the implied
ones, so a model asks reasonable questions but rarely the complete set
from the claim alone \cite{chen2022claimdecomp}. Varifocal generates
questions from several focal points of the claim and its metadata, such
as its speaker and date, without presuming that the answer is already in
a supplied passage \cite{ousidhoum2022varifocal}; FACTIFY-5WQA builds its
questions from the claim's \emph{who, what, when, where} and \emph{why}
slots, found by semantic role labelling \cite{rani2023factify5wqa}.
AVeriTeC adopts the question as the unit of evidence: its annotators
record each claim's evidence as question--answer pairs, and a system's
evidence is scored by matching its pairs to theirs
\cite{schlichtkrull2023averitec}. The systems that led its 2024 shared
task \cite{schlichtkrull2024sharedtask} all ask before they search:
InFact poses ten questions and lets the model write the search queries
for each, which its authors found better than searching the question
literally \cite{rothermel2024infact}; HerO, built on open models only,
generates the questions with an 8B Llama~3 prompted with ten examples
\cite{yoon2024hero}.

\paragraph{Splitting the claim.} Other pipelines split a claim into
simpler sub-claims and verify each: ProgramFC as the steps of a program
\cite{pan2023programfc}, FActScore as atomic facts
\cite{min2023factscore}, Factcheck-Bench as claims that each hold one
property or fact \cite{wang2024factcheckbench}. Chen et al.\ add the
constraint this work also faces, evidence published before the claim on
the open web, and find that retrieved evidence usually covers some
aspects of a claim and rarely all of them \cite{chen2024wild}. Whether
splitting helps is contested. FActScore's result depends on the
decomposition method chosen \cite{wanner2024closer}. Hu, Long and Wang
find a trade-off: simpler sub-claims help a weaker verifier, but every
sub-claim adds noise, and with a strong verifier the noise can outweigh
the gain; more sub-claims help at first and then hurt
\cite{hu2025decomposition}.

\paragraph{Making the claim stand alone.} A sentence taken from an
article leans on its context: ``the agency'', ``that year'', a subject
named only in the headline. Decontextualisation rewrites it so that it
can be read on its own \cite{choi2021decontextualization}. Gunjal and
Durrett show the cost of going too far the other way: atomic facts can
lose the context that makes them checkable, and they propose
\emph{molecular} facts, which carry the least added context that makes
them unambiguous \cite{gunjal2024molecular}.

\paragraph{Small models.} The step need not take a large model.
ClaimCheck plans its web queries with a 4B model and reports 76.4\%
accuracy on AVeriTeC's development split; turning the model's reasoning
off in the planning step alone cost about four points in its ablation
\cite{putta2025claimcheck}.

\paragraph{This work.} By default the pipeline has no such step: every
query is built from the sentence's own words
(Section~\ref{sec:design-fact-checking}). The one built and measured so
far takes no model. It reads the claim's \emph{checkable point} from the
sentence by rule --- figures with their units, dates, multi-word names,
the words of the finding before those of its framing --- and restores
the article's subject when the sentence leaves it out. It is close to
the literal subquestions and the 5W slots above, and to
decontextualisation, but it writes no question and names no kind of
source. It is kept as it is, as the baseline against which a
question-based plan is to be measured
(Sections~\ref{sec:results-pre-verdict} and~\ref{subsec:future-further}).
